\section{Applicability and Source-Bound Sensitivity}\label{sec:r36-review}
\subsection{Two adaptation questions and heterogeneous costs}
R4/R5 refer to protected-data fitting and a matched training-fit contrast. For a non-learning attack the historical ``--'' describes that specific contrast only; an adaptive strategy that adjusts probes, auxiliary information or operating conditions can still be relevant. The broadened sensitivity projection converts the historical ``--'' cells to unknown (?), not to U or a verified missing experiment. The historical record labels remain unchanged, and a separate source review would be needed to recode broad adaptive evaluation. Table~\ref{tab:r36-sensitivity} reports the resulting eligible denominators.
\begin{table}[t]\centering\footnotesize
\caption{Two fit-applicability and cost-placement interpretations over the same 31 privacy/attribution records. Entries are bound/eligible records. Broadening maps fit-inapplicable records to unknown, not to a verified absence of adaptation.}\label{tab:r36-sensitivity}
\begin{tabularx}{\columnwidth}{@{}Yrr@{}}\toprule
Query & Fit scope & Broader unknown\\\midrule
Matched adaptation + inference & 2/15 & 2/22\\
All five bindings & 0/13 & 0/19\\
\midrule
Inference + cost & 5/21 & 4/21$^{\dagger}$\\
\bottomrule\end{tabularx}
\par\smallskip\raggedright $^{\dagger}$ Requires explicit R3 placement as well as R7/R9, or treats SPEC's payload-only cost as P for a wire-level query. The same four records remain, with distinct cost objects. Source cells and original query labels remain unchanged.
\end{table}

 The full item memberships and hashes are in \path{data/r36-sensitivity.json}.

SPEC's REST padding record reports an inference result and a payload-size ratio at its evaluated configuration (inspected v4, Section 6.1 and Appendix D Table 2; \url{https://arxiv.org/html/2411.01076v4}). It receives R9 B for that stated payload object, while R3 is P because placement/execution binding is incomplete. It does not measure wire bytes or end-to-end cost. Requiring R3 B, or treating payload-only cost as partial for a wire-cost question, removes SPEC from the inference-plus-cost positives; the four remaining records still refer to distinct cost objects. The five original positives are not a common cost ranking.

A leave-one-family-out computation reruns each query over the 19 Type I/II mitigation families. For the triple query the two CachePrune record positives both disappear when CachePrune is omitted, and remain for each of the other 18 omissions. They represent two target outcomes in one configuration. Removing all final-search additions gives 0/15 eligible triple records; all queries retain their own declared eligibility denominators. The projection preserves the frozen source cells and does not establish that unreviewed evidence is absent.

\subsection{Search frame and boundary cases}
The 615 occurrences and 565 deduplicated records in Section~\ref{app:search-update} are a completed, finite retrieval update, not a PRISMA census of all venues. Its disposition table and query strings give the reproducible route; 392 records were routed to the two stated exclusions and the remaining 173 to the other reported categories. The initial seeds and venue screening overlap that update and must not be added as independent PRISMA stages. Hardware and architecture venues including ISCA, MICRO, HPCA, ASPLOS and CHES/TCHES were not comprehensively screened. The October 21 closure is a snapshot, not a continuously updated claim.

\begin{figure}[ht]\centering
\begin{tikzpicture}[node distance=4mm, every node/.style={font=\footnotesize,align=center}, box/.style={draw,rounded corners,minimum width=54mm,minimum height=7mm}]
\node[box] (a) {Completed update: 615 occurrences};
\node[box,below=of a] (b) {565 unique records};
\node[box,below left=10mm and 2mm of b] (c) {392 excluded: target or scope};
\node[box,below right=10mm and 2mm of b] (d) {173 routed to other categories};
\draw[-{Latex}] (a)--node[right]{50 duplicates merged}(b);
\draw[-{Latex}] (b)--(c); \draw[-{Latex}] (b)--(d);
\end{tikzpicture}
\caption{Finite retrieval update only. Other routes include already indexed studies, new and adjacent records, background and unresolved availability; full counts appear in Table~\ref{tab:finite-flow}. Overlapping seed and venue routes are excluded.}\label{fig:r36-search-flow}
\end{figure}

Three adjacent decisions illustrate the operational boundary. CacheSolidarity v1 (\url{https://arxiv.org/html/2603.10726v1}) has an auxiliary timing/cache target and a candidate-prefix experiment in Section 5.2.3, Figure 9. The archived family index incorrectly points to Figure 7, which is a benign-workload figure; Figure 9 is the corrected locator. The inspected v1 presents attack traces and qualitative interpretation; whether that meets C5 when a scored inference result is required is a boundary decision, not proof of no attack. Inclusion under a broader C5 requires a separate re-audit and count update. TOPR v4 (\url{https://arxiv.org/html/2512.16310v4}) studies information exposed through response/reasoning content; under the current semantic-only exclusion its inspected analysis does not isolate the auxiliary-only C3 condition. SecEmb (\url{https://bearhw.ece.vt.edu/content/dam/bearhw_ece_vt_edu/publications/2025_HPCA_SecEmb.pdf}) pairs its attack demonstration with DLRM embedding operations; its GPT-2 measurements concern protection cost/quality rather than a paired GPT-2 attacker recovery result. These are scoped selection explanations, not claims that the studies lack security value. Other adjacent defense decisions remain open. SICC (\url{https://arxiv.org/html/2609.26076v1}, Sections 3 and 5) studies private-state-dependent message form in addition to content; DP-MemView (\url{https://arxiv.org/html/2608.03130v1}, Appendix B.5) includes an internal-score side-channel diagnostic. Whether their evaluated receiver and isolated signal satisfy C3--C5 requires a record-level adjudication, not a blanket semantic-only exclusion. PlanTwin's re-identification experiment (\url{https://arxiv.org/html/2603.18377v2}, Section 5.6) uses disclosed schema fields; its relation to the auxiliary-only C3 boundary also needs explicit adjudication. PriCon's accessible HTML (\url{https://arxiv.org/html/2609.32642v1}) supplies only an abstract; its detailed eligibility remains unresolved. No complete adjacent-defense exclusion ledger is claimed.
